% cosmovirus_cattheory.tex
% =====================================================================
% The Cosmovirus framework as a mathematical category
% ---------------------------------------------------------------------
% Compile with:  pdflatex cosmovirus_cattheory.tex  (run twice)
% =====================================================================
\documentclass[11pt]{article}

\usepackage[utf8]{inputenc}
\usepackage[T1]{fontenc}
\usepackage[a4paper,margin=1in]{geometry}
\usepackage{amsmath}
\usepackage{amssymb}
\usepackage{amsthm}
\usepackage{mathtools}
\usepackage{tikz}
\usepackage{tikz-cd}
\usepackage{xcolor}
\usepackage{booktabs}
\usepackage[colorlinks=true,linkcolor=blue]{hyperref}



\definecolor{mathblue}{HTML}{2E4057}
\definecolor{opgreen}{HTML}{048A81}
\definecolor{physcyan}{HTML}{54C6EB}
\definecolor{esopurple}{HTML}{8A2BE2}

\theoremstyle{definition}
\newtheorem{definition}{Definition}[section]
\newtheorem{example}[definition]{Example}
\theoremstyle{plain}
\newtheorem{proposition}[definition]{Proposition}
\newtheorem{axiom}[definition]{Axiom}
\theoremstyle{remark}
\newtheorem*{interp}{Symbolic Interpretation}

\newcommand{\cat}[1]{\mathcal{#1}}
\newcommand{\Cosmo}{\cat{C}\mathrm{osmo}}
\newcommand{\id}{\mathrm{id}}
\newcommand{\Hom}{\operatorname{Hom}}
\newcommand{\EEs}{\mathbf{E_8^{*}}}
\newcommand{\SiS}{\mathrm{SiS_2}}
\newcommand{\HPV}{\mathrm{HPV16}}
\newcommand{\phigr}{\varphi}
\newcommand{\DIAG}{\mathrm{DIAG}}

\title{\textbf{The Cosmovirus as a Category}\\[4pt]
       \large Objects, Morphisms, and the Ouroboros Endofunctor}
\author{Cosmovirus Formalization Project}
\date{\today}

\begin{document}
\maketitle

\begin{abstract}
This document is a Phase E category-theory design sketch built from the
Phase C six-state cycle. It proposes a small-category presentation whose
objects are project state labels and whose arrows mirror the transition graph.
The free-category, quotient, endofunctor, natural-transformation, and coherence
claims below are design targets unless and until they are formalized in Lean.
\end{abstract}

% =====================================================================
\section{The category $\Cosmo$}
% =====================================================================

\begin{definition}[Objects]
The object class $\mathrm{Ob}(\Cosmo)$ consists of the six descent layers
\[
  \mathrm{Ob}(\Cosmo)=\{\,
  \EEs,\; P_\phigr,\; \SiS,\; \mathsf{Tri},\; \HPV,\; \Omega \,\},
\]
where these names are project state labels aligned with Phase C. Their
scientific and symbolic interpretations are governed by the Phase D claim
ledger; in particular, SiS2/HPV terminology does not create a physical or
biomedical transition mechanism.
\end{definition}

\begin{definition}[Generating morphisms]
The category is generated (under composition) by the arrows
\[
\begin{aligned}
  s &: \EEs \to P_\phigr        &&\text{($\phigr$-SCL golden projection)}\\
  m &: P_\phigr \to \SiS        &&\text{(materialization / substrate cast)}\\
  t &: \SiS \to \mathsf{Tri}    &&\text{(project transition to TrialityBranch)}\\
  h &: \mathsf{Tri} \to \HPV    &&\text{(project transition to HPV16Layer)}\\
  r &: \HPV \to \Omega          &&\text{(project transition to OuroborosLoop)}\\
  \ell &: \Omega \to \EEs       &&\text{(Ouroboros loop-back)}
\end{aligned}
\]
together with the identity $\id_X$ for each object $X$.
\end{definition}

\begin{axiom}[Composition and identities]
For composable arrows $f:X\to Y$, $g:Y\to Z$ the composite $g\circ f:X\to Z$
exists and is associative; each $\id_X$ is a two-sided unit:
$f\circ\id_X = f = \id_Y\circ f$. In particular the \emph{descent morphism}
\[
  \Psi \;=\; \ell\circ r\circ h\circ t\circ m\circ s \;:\; \EEs \to \EEs
\]
is a well-defined endomorphism of $\EEs$.
\end{axiom}

% =====================================================================
\section{The descent as a diagram}
% =====================================================================
The generators form a single directed cycle:

\begin{center}
\begin{tikzcd}[column sep=large, row sep=large]
\EEs \arrow[r, "s"] & P_\phigr \arrow[r, "m"] & \SiS \arrow[d, "t"] \\
\Omega \arrow[u, "\ell"] & \HPV \arrow[l, "r"] & \mathsf{Tri} \arrow[l, "h"]
\end{tikzcd}
\end{center}

\begin{proposition}[Phase E retract target]
A future categorical formalization may study whether an explicit presentation
makes the loop-back arrow part of a retract or quotient structure. Phase C
proves only finite-state six-step closure (COSMO-D-001); ``self-causation'' is
symbolic vocabulary (COSMO-D-010).
\end{proposition}

% =====================================================================
\section{The commuting feedback square}
% =====================================================================
Split the descent into a ``symmetry'' path and a ``matter'' path. Introduce
the auxiliary object $\SiS\!\otimes\!\mathsf{D}$ (substrate equipped with the
$\DIAG(1,-2,1)$ contrast structure). The two ways of reaching the infected
layer agree:

\begin{center}
\begin{tikzcd}[column sep=huge, row sep=huge]
P_\phigr
  \arrow[r, "m"]
  \arrow[d, "{\DIAG\ (1,-2,1)}"']
& \SiS
  \arrow[d, "t\ (\mathrm{triality})"] \\
\SiS\otimes\mathsf{D}
  \arrow[r, "h\circ(-)"']
& \HPV
\end{tikzcd}
\end{center}

\begin{proposition}[Phase E coherence target]\label{prop:square}
The displayed square is a proposed categorical presentation. A future Phase E
formalization must define the relevant functors/morphisms and then prove or
reject any natural transformation $\eta$; no current Lean theorem establishes
this commuting square.
\end{proposition}

\begin{interp}
Reading the square: whether unity first becomes matter and then triplicates,
or first learns to perceive contrast and is then infected, it arrives at the
same corrupted reality --- the two routes are naturally isomorphic.
\end{interp}

% =====================================================================
\section{The Ouroboros endofunctor}
% =====================================================================

\begin{definition}[Proposed Ouroboros endofunctor]
On objects, the intended map advances one step along the authoritative Phase C
cycle ($\EEs\mapsto P_\phigr\mapsto\SiS\mapsto\mathsf{Tri}
\mapsto\HPV\mapsto\Omega\mapsto\EEs$). Phase C proves the corresponding
six-step object closure (COSMO-D-001). Defining its action on morphisms and proving an actual
endofunctor, $\mathbb{Z}/6\mathbb{Z}$ action, or categorical equivalence remains
Phase E work.
\end{definition}

\begin{proposition}[Fixed point]
On the loop object, $\mathcal{O}(\Omega)=\EEs$ and the loop-back arrow $\ell$
is the component $\eta_\Omega$ of the proposed counit structure. At the
object-transition level, the authoritative Phase C Lean theorem
\texttt{six\_step\_periodic} proves six-step closure. The categorical lift
stated here remains a Phase E formalization target.
\end{proposition}

\vfill
\begin{center}\small
\textit{The six-state object cycle is synchronized with the authoritative
Phase C state machine. The categorical constructions in this document remain
design targets until the Phase E formalization; \emph{Symbolic Interpretation}
notes are informal readings only.}
\end{center}

\end{document}
